\documentclass[11pt]{article}

\usepackage[T1]{fontenc}
\usepackage[utf8]{inputenc}
\usepackage[a4paper,margin=1in]{geometry}
\usepackage{amsmath,amssymb,amsthm}
\usepackage{mathtools}
\usepackage{hyperref}

\newtheorem{theorem}{Theorem}
\newtheorem{corollary}{Corollary}
\newtheorem{remark}{Remark}
\newtheorem{example}{Example}
\newtheorem{assumption}{Assumption}

\newcommand{\R}{\mathbb{R}}
\newcommand{\safeSet}{\mathcal{C}}
\newcommand{\Dt}{\frac{d}{dt}}
\newcommand{\Trec}{T_{\mathrm{rec}}}
\newcommand{\taumax}{\tau_{\max}}
\newcommand{\tauC}{\tau_{\mathcal C}}

\title{Regional Finite-Time Recovery for Higher-Order Control Barrier Functions}

\author{Huatian Zhu \and GPT-5.5}

\date{June 15, 2026}

\begin{document}

\maketitle

\begin{abstract}
This note studies finite-time recovery from unsafe states for safety constraints with arbitrary relative degree. For a control-affine nonlinear system, we consider a smooth higher-order barrier inequality generated by a repeated constant-gain differential operator. The resulting constraint is affine in the control input and is therefore compatible with quadratic-program-based safety filters. A general regional finite-time recovery condition is derived through a Taylor-type lower bound after an exponential change of variables. A simpler sufficient condition is also given, requiring positivity of the first auxiliary variable and nonnegativity of the remaining lower-order auxiliary variables. The result preserves smoothness and provides an explicit recovery-time bound on a regional subset of unsafe initial states, under the assumption that the barrier inequality is feasible and enforced before recovery.
\end{abstract}

\section{Introduction}

Control barrier functions provide a systematic method for enforcing safety constraints in nonlinear control systems. For control-affine systems, barrier inequalities that are affine in the control input can be incorporated directly into quadratic programs, yielding safety filters that minimally modify a nominal controller while enforcing prescribed safety conditions \cite{ames2014control,ames2019control}.

When the safety constraint has relative degree one, the control input appears directly in the first derivative of the barrier function. For higher-relative-degree constraints, higher-order control barrier function constructions introduce auxiliary barrier variables until the control input appears. Related approaches include exponential control barrier functions and high-order control barrier functions \cite{nguyen2016exponential,xiao2019control}. These constructions are commonly used for forward invariance of a safe set.

In many applications, the system may start outside the safe set. In this case, forward invariance alone does not address the full safety objective. A distinct question is whether the imposed barrier condition can recover safety in finite time. This note studies such a recovery question for higher-order control barrier functions.

The main contribution is a regional finite-time recovery result for arbitrary relative degree. The recovery condition is expressed through a polynomial determined by the initial auxiliary barrier variables. A simple sufficient region is then obtained by imposing positivity of the first auxiliary variable and nonnegativity of the remaining lower-order auxiliary barrier variables. The construction remains smooth and affine in the control input.

\section{Preliminaries and Problem Formulation}

Consider a control-affine nonlinear system
\begin{equation}
\dot x = f(x)+g(x)u,
\label{eq:system}
\end{equation}
where $x\in\R^n$, $u\in\R^m$, and $f,g$ are sufficiently smooth. Let
\begin{equation}
\safeSet = \{x\in\R^n : h(x)\ge0\}
\label{eq:safe-set}
\end{equation}
be the desired safe set, where $h:\R^n\to\R$ is sufficiently smooth.

We assume that $h$ has relative degree $r\ge2$ with respect to \eqref{eq:system} on the region of interest. That is,
\begin{equation}
L_gL_f^j h(x)=0,
\qquad j=0,\dots,r-2,
\label{eq:relative-degree-zero}
\end{equation}
and
\begin{equation}
L_gL_f^{r-1}h(x)\neq0.
\label{eq:relative-degree-nonzero}
\end{equation}
Thus, along sufficiently regular system trajectories,
\begin{equation}
h^{(j)}(t)=L_f^j h(x(t)),
\qquad j=0,\dots,r-1,
\end{equation}
and
\begin{equation}
h^{(r)}(t)
=
L_f^r h(x(t))
+
L_gL_f^{r-1}h(x(t))u(t).
\end{equation}

The objective is to derive sufficient conditions under which a trajectory starting from an unsafe state,
\begin{equation}
h(x_0)<0,
\label{eq:unsafe-initial-state}
\end{equation}
enters $\safeSet$ in finite time while satisfying a smooth higher-order barrier constraint before recovery.

\begin{remark}[Regularity convention]
The statements below are written in a form that accommodates both classical smooth controls and the piecewise-smooth or measurable controls often generated by optimization-based filters. On every compact time interval under consideration, we assume that the trajectory exists and that the scalar function
\begin{equation}
w(t)=e^{kt}h(x(t))
\end{equation}
is $(r-1)$ times continuously differentiable, with $w^{(r-1)}$ absolutely continuous. Then Taylor's formula with integral remainder is valid, and inequalities involving $w^{(r)}$ may be interpreted almost everywhere. In the classical case, this reduces to the usual pointwise derivative interpretation.
\end{remark}

\section{Main Results}

\subsection{Higher-Order Barrier Construction}

Fix $k>0$. Define the auxiliary barrier variables
\begin{equation}
\psi_0 = h,
\end{equation}
and, for $i=1,\dots,r$,
\begin{equation}
\psi_i
=
\left(\Dt+k\right)^i h.
\label{eq:psi-definition}
\end{equation}
Equivalently,
\begin{equation}
\psi_i
=
\dot\psi_{i-1}+k\psi_{i-1},
\qquad i=1,\dots,r.
\label{eq:recursive-psi}
\end{equation}

For $i=0,\dots,r-1$, the relative-degree assumption implies that $\psi_i$ is independent of the control input and can be written as
\begin{equation}
\psi_i(x)
=
\sum_{j=0}^{i}
\binom{i}{j}
k^{i-j}L_f^j h(x).
\label{eq:psi-expansion}
\end{equation}

The higher-order barrier recovery constraint is
\begin{equation}
\psi_r(x,u)\ge0.
\label{eq:hocbf-constraint}
\end{equation}
Expanding \eqref{eq:hocbf-constraint} gives
\begin{equation}
L_f^r h(x)
+
L_gL_f^{r-1}h(x)u
+
\sum_{j=0}^{r-1}
\binom{r}{j}
k^{r-j}L_f^j h(x)
\ge0.
\label{eq:expanded-hocbf}
\end{equation}
This constraint is affine in $u$, and hence can be incorporated directly into a quadratic-program-based safety filter.

\subsection{Regional Finite-Time Recovery}

We now present the main recovery result. The condition is expressed in terms of the initial values of the auxiliary barrier variables.

\begin{theorem}[Regional finite-time recovery]
\label{thm:general-recovery}
Let $x(t)$ be a solution of \eqref{eq:system} on a maximal interval $[0,\taumax)$, where $\taumax\in(0,\infty]$. Define the first entrance time into the safe set by
\begin{equation}
\tauC
=
\inf\{t\in[0,\taumax): h(x(t))\ge0\},
\label{eq:first-entrance-time}
\end{equation}
with the convention that $\inf\emptyset=\infty$.

Suppose that $h(x_0)<0$ and that, for every compact interval $[0,T]\subset[0,\taumax)$, the function
\begin{equation}
w(t)=e^{kt}h(x(t))
\end{equation}
is $(r-1)$ times continuously differentiable and $w^{(r-1)}$ is absolutely continuous.

Assume that, for almost every $t<\tauC$,
\begin{equation}
\psi_r(x(t),u(t))\ge0.
\label{eq:enforced-constraint}
\end{equation}
Let
\begin{equation}
P_{r-1}(t)
=
\sum_{i=0}^{r-1}
\frac{\psi_i(x_0)}{i!}t^i,
\label{eq:recovery-polynomial}
\end{equation}
where $\psi_0(x_0)=h(x_0)$.

If the set
\begin{equation}
S
=
\{T\in(0,\taumax): P_{r-1}(T)\ge0\}
\label{eq:crossing-set}
\end{equation}
is nonempty, then the trajectory enters the safe set $\safeSet$ in finite time. Moreover,
\begin{equation}
\Trec
=
\tauC
\le
\inf S
=
\inf
\{
T\in(0,\taumax):
P_{r-1}(T)\ge0
\}.
\label{eq:general-time-bound}
\end{equation}
\end{theorem}

\begin{proof}
Define
\begin{equation}
w(t)=e^{kt}h(x(t)).
\label{eq:w-definition}
\end{equation}
By direct differentiation,
\begin{equation}
w^{(i)}(t)
=
e^{kt}\psi_i(x(t)),
\qquad i=0,\dots,r,
\label{eq:w-psi-relation}
\end{equation}
where the identity for $i=r$ is understood almost everywhere under the stated regularity assumptions. In particular, since $e^{kt}>0$, the imposed constraint \eqref{eq:enforced-constraint} implies
\begin{equation}
w^{(r)}(t)\ge0
\label{eq:w-r-nonnegative}
\end{equation}
for almost every $t<\tauC$.

Let $T\in S$ be arbitrary. If $\tauC\le T$, then the desired bound holds for this $T$.

Now suppose, for contradiction, that $\tauC>T$. Then $h(x(t))<0$ for all $t\in[0,T]$, and therefore the barrier constraint is enforced almost everywhere on $[0,T]$. Hence $w^{(r)}(t)\ge0$ for almost every $t\in[0,T]$.

By Taylor's formula with integral remainder,
\begin{equation}
w(T)
=
\sum_{i=0}^{r-1}
\frac{w^{(i)}(0)}{i!}T^i
+
\int_0^T
\frac{(T-\tau)^{r-1}}{(r-1)!}
w^{(r)}(\tau)\,d\tau.
\label{eq:taylor-integral}
\end{equation}
Using \eqref{eq:w-psi-relation} at $t=0$, we have
\begin{equation}
w^{(i)}(0)=\psi_i(x_0),
\qquad i=0,\dots,r-1.
\end{equation}
Since $w^{(r)}(\tau)\ge0$ almost everywhere on $[0,T]$, it follows that
\begin{equation}
e^{kT}h(x(T))
=
w(T)
\ge
\sum_{i=0}^{r-1}
\frac{\psi_i(x_0)}{i!}T^i
=
P_{r-1}(T).
\label{eq:taylor-lower-bound}
\end{equation}
Because $T\in S$, we have $P_{r-1}(T)\ge0$. Therefore,
\begin{equation}
e^{kT}h(x(T))\ge0.
\end{equation}
Since $e^{kT}>0$,
\begin{equation}
h(x(T))\ge0.
\end{equation}
This contradicts $\tauC>T$. Hence $\tauC\le T$.

Because the above argument holds for every $T\in S$, we obtain
\begin{equation}
\tauC\le\inf S.
\end{equation}
Thus the trajectory enters $\safeSet$ in finite time, and the stated recovery-time bound follows.
\end{proof}

\begin{remark}[The second-order case]
When $r=2$, the recovery polynomial reduces to
\begin{equation}
P_1(t)
=
h(x_0)+\psi_1(x_0)t.
\end{equation}
Therefore, under $h(x_0)<0$, the polynomial crossing condition in Theorem~\ref{thm:general-recovery} is equivalent to
\begin{equation}
\psi_1(x_0)>0.
\end{equation}
The corresponding recovery-time bound becomes
\begin{equation}
\Trec
\le
\frac{-h(x_0)}{\psi_1(x_0)}.
\end{equation}
Thus, the second-order recovery statement is obtained as a direct consequence of the general theorem rather than as an independent assumption.
\end{remark}

\subsection{A Simple Sufficient Recovery Region}

The polynomial crossing condition in Theorem~\ref{thm:general-recovery} is general but may be inconvenient to check. The following corollary provides a simpler sufficient condition.

\begin{corollary}[Simple regional recovery condition]
\label{cor:simple-region}
Suppose that the hypotheses of Theorem~\ref{thm:general-recovery} hold. If
\begin{equation}
h(x_0)<0,
\label{eq:simple-h-negative}
\end{equation}
\begin{equation}
\psi_1(x_0)>0,
\label{eq:simple-psi-one-positive}
\end{equation}
and
\begin{equation}
\psi_i(x_0)\ge0,
\qquad i=2,\dots,r-1,
\label{eq:simple-higher-nonnegative}
\end{equation}
then the trajectory enters $\safeSet$ in finite time. Moreover,
\begin{equation}
\Trec
\le
\frac{-h(x_0)}{\psi_1(x_0)},
\label{eq:simple-time-bound}
\end{equation}
provided that
\begin{equation}
\frac{-h(x_0)}{\psi_1(x_0)}<\taumax.
\label{eq:simple-time-existence}
\end{equation}
\end{corollary}

\begin{proof}
Under \eqref{eq:simple-psi-one-positive} and \eqref{eq:simple-higher-nonnegative}, the recovery polynomial satisfies
\begin{equation}
P_{r-1}(t)
=
h(x_0)
+
\psi_1(x_0)t
+
\sum_{i=2}^{r-1}
\frac{\psi_i(x_0)}{i!}t^i
\ge
h(x_0)+\psi_1(x_0)t
\end{equation}
for all $t\ge0$. Therefore, at
\begin{equation}
T=
\frac{-h(x_0)}{\psi_1(x_0)},
\end{equation}
we have
\begin{equation}
P_{r-1}(T)\ge0.
\end{equation}
By \eqref{eq:simple-time-existence}, this $T$ belongs to $(0,\taumax)$. Hence $T\in S$, and the result follows from Theorem~\ref{thm:general-recovery}.
\end{proof}

\begin{remark}[Why the polynomial condition is needed]
The condition $\psi_1(x_0)>0$ only states that the first derivative of the exponentially weighted barrier is initially positive. For $r\ge3$, this does not control the remaining Taylor coefficients of $e^{kt}h(x(t))$.

For example, consider
\begin{equation}
w(t)=-1+t-\frac{t^2}{2}.
\end{equation}
Then
\begin{equation}
w(0)<0,
\qquad
w'(0)>0,
\qquad
w^{(3)}(t)=0,
\end{equation}
but
\begin{equation}
w(t)<0
\end{equation}
for all $t\ge0$. Taking $k=1$ and $h(t)=e^{-t}w(t)$ gives
\begin{equation}
\left(\Dt+1\right)^3h(t)=0
\end{equation}
and
\begin{equation}
\psi_1(0)>0,
\end{equation}
while $h(t)<0$ for all $t\ge0$.

Thus, for higher-order constraints, positivity of the first auxiliary variable alone is insufficient. Some information about the remaining lower-order auxiliary variables, or equivalently the polynomial crossing condition in Theorem~\ref{thm:general-recovery}, is needed. This example is a scalar comparison-level counterexample; it demonstrates the obstruction in the Taylor comparison argument itself.
\end{remark}

\section{Discussion and Conclusion}

The result is regional because it does not claim recovery from every unsafe state. The admissible unsafe initial states are those for which the initial auxiliary variables generate a recovery polynomial that crosses zero in positive time before the trajectory ceases to exist. A simple sufficient region is
\begin{equation}
\{
x\in\R^n:
h(x)<0,\
\psi_1(x)>0,\
\psi_i(x)\ge0,\ i=2,\dots,r-1
\}.
\label{eq:simple-region-set}
\end{equation}
For initial states in this region, the linear part of the recovery polynomial already crosses zero, and the higher-order Taylor coefficients do not oppose the comparison argument.

The proposed construction preserves several useful features. It is smooth, uses constant gains, and produces a constraint that is affine in the control input. These properties are compatible with optimization-based safety filters. The analysis assumes that the barrier constraint is feasible and enforced before recovery; feasibility of the resulting quadratic program is a separate issue and depends on the control authority and any additional input constraints.

The result should also be distinguished from a forward-invariance guarantee. The theorem proves finite-time entrance into $\safeSet$ from a regional subset of unsafe initial states under the imposed higher-order barrier inequality before recovery. Maintaining safety after entrance may require continuing to enforce an appropriate invariance condition or combining the recovery mechanism with a standard control barrier function condition inside or near the safe set.

In summary, smooth higher-order barrier inequalities can provide finite-time recovery from unsafe states for arbitrary relative degree $r\ge2$, but the recovery guarantee is intrinsically regional. The result requires either a polynomial crossing condition or a simple set of nonnegativity conditions on the lower-order auxiliary barrier variables, together with trajectory existence and the regularity needed for the Taylor comparison argument.

\begin{thebibliography}{9}

\bibitem{ames2014control}
A. D. Ames, J. W. Grizzle, and P. Tabuada,
``Control barrier function based quadratic programs with application to adaptive cruise control,''
in \emph{Proceedings of the 53rd IEEE Conference on Decision and Control},
2014, pp. 6271--6278.

\bibitem{nguyen2016exponential}
Q. Nguyen and K. Sreenath,
``Exponential control barrier functions for enforcing high relative-degree safety-critical constraints,''
in \emph{Proceedings of the American Control Conference},
2016, pp. 322--328.

\bibitem{xiao2019control}
W. Xiao and C. Belta,
``Control barrier functions for systems with high relative degree,''
in \emph{Proceedings of the 58th IEEE Conference on Decision and Control},
2019, pp. 474--479.

\bibitem{ames2019control}
A. D. Ames, S. Coogan, M. Egerstedt, G. Notomista, K. Sreenath, and P. Tabuada,
``Control barrier functions: Theory and applications,''
in \emph{Proceedings of the 18th European Control Conference},
2019, pp. 3420--3431.

\end{thebibliography}

\end{document}
